The observability practices of monitoring, tracing and logging play an important role in ensuring the reliability of cloud-native microservice architectures. 
Microservices provide many advantages over their monolithic predecessors~\cite{Fowler_Microservices_2014}, however the resulting applications can be difficult to troubleshoot when a fault occurs\cite{Niedermaier_ObservabilityInterviewStudy_2019,Zhang_MicroserviceSurvey_ThresholdingHard_2019}. This is because each request often travels through multiple independently developed services, deployed on an intricate and evolving stack of software platforms. To address this challenge, observability systems have evolved accordingly, %
now consisting of distributed services that collect observability data from all levels of runtime platforms as well as from the application level through built-in and developer-defined instrumentation.
Hence, using these observability systems requires a multitude of hard decisions in terms of instrumentation points, observability services and configuration. 
When employed correctly, observability can help developers identify faults quickly, but improper settings can also obfuscate faults \cite{Zhang_MicroserviceSurvey_ThresholdingHard_2019,Borges_TracingServerless_2021} and increase latencies and cost.

Still, designing %
the observability of microservice applications is not trivial. It implies often overlooked tasks like the tool-dependent configuration of parameters, setting alerts and adding custom instrumentation code. These decisions need to be weighed against observability-related trade-offs, e.g., the performance overhead on the application, plus the cost overhead associated with operating the observability infrastructure.   %
However, these architectural design decisions are not supported by systematic methods. Instead, decisions are made based on previous experience and “professional intuition” of the practitioners \cite{Zhang_MicroserviceSurvey_ThresholdingHard_2019,Vale_MicroserviceQualityMeasurementIndustrySurvey_2022}, or developers may end up eagerly instrumenting and changing configuration parameters in a reactive manner after a problem occurs \cite{Niedermaier_ObservabilityInterviewStudy_2019}. To weigh between different design alternatives and arrive at an appropriate configuration that justifies the effort, practitioners must have a method to assess the effectiveness of their observability. 

Consequently, the need to evaluate the degree of observability of a system has been recognized by research and industry alike \cite{Tamburri_MVCObservability_2018,Google_MonitoringMeasurement_2023,Ahmed_EffectivenessOfAPM_2016,Vale_MicroserviceQualityMeasurementIndustrySurvey_2022}. With suitable metrics, it would be possible to obtain a concrete understanding of the quality of the observability, which could then be tracked and continuously improved over time. However, observability has been challenging to quantify so far \cite{Tamburri_MVCObservability_2018}. Most research has focused either on qualitative assessments of observability tooling \cite{Janes_TracingQualitativeAnalysis_2022}, or solely looks at cost and overhead  \cite{Ernst_OfflineTraceGeneration_2021,Reichelt_OverheadOpenTelemetryEtc,Dinga_EnergyEfficiencyOfMonitoring_2023}, failing to address the question of \textit{effectiveness} of observability. 
Currently, there are no mechanisms for practitioners to quantify or compare the observability of their applications.

In this paper, we argue for a systematic, reproducible and comparative assessment of observability design decisions. %
Our focus lies on the observability of faults specifically. Similar to software quality measures like test coverage, we aim to make fault observability a testable and quantifiable system property. This, in turn, would help developers identify unobservable faults in their application, thus guiding the configuration and instrumentation process. 

Towards this goal, we present the following contributions: \begin{enumerate}
    \item A model for understanding the scale and scope of observability design decisions.
    \item The concept for testable and quantifiable fault observability metrics.%
    \item Design, implementation and demonstration of \name{}, a tool to automate the process of observability assessments.
\end{enumerate}
The paper is structured as follows: Section \ref{sec:background}  discusses related work.  Section \ref{sec:model1}  provides background and models the observability design space of cloud-native architectures. Section \ref{sec:concept} presents our observability metrics and experiment method. Section \ref{sec:architecture} introduces our supporting experiment tool \name{}. Section \ref{sec:eval} evaluates our method and tool through exemplary experiments. Section \ref{sec:discussion} discusses suitability and limitations. Section \ref{sec:conclusion} concludes with future work.
